\section{Classi Definite}
\label{sec:classi_definite}

\subsection{Panoramica}

Le classi definite (\textit{defined classes}) costituiscono uno degli aspetti
pi\`u avanzati dell'ontologia. A differenza delle classi primitive, che richiedono
l'asserzione esplicita \texttt{rdf:type}, le classi definite specificano
condizioni necessarie e sufficienti (\texttt{owl:equivalentClass}) che permettono
a un reasoner di \textbf{inferire automaticamente} l'appartenenza di un'istanza.

L'ontologia definisce \textbf{17 classi definite}, raggruppabili in sei categorie
logiche. La Tabella~\ref{tab:definite} le elenca con la sintassi DL e il tipo
di condizione.

\begin{longtable}{lllp{4cm}}
\caption{Classi definite}\label{tab:definite}\\
\toprule
\textbf{Classe} & \textbf{Tipo} & \textbf{Sintassi DL} & \textbf{Costrutti OWL 2} \\
\midrule
\endfirsthead
\multicolumn{4}{c}{\tablename\ \thetable{} -- continua} \\
\toprule
\textbf{Classe} & \textbf{Tipo} & \textbf{Sintassi DL} & \textbf{Costrutti OWL 2} \\
\midrule
\endhead
\bottomrule
\endlastfoot

% Mode-derived
MultiplayerGame     & S/N & VideoGame $\sqcap$ $\exists$hasGameMode.MultiplayerMode   & intersectionOf, someValuesFrom \\
SinglePlayerGame    & S/N & VideoGame $\sqcap$ $\exists$hasGameMode.SinglePlayerMode  & intersectionOf, someValuesFrom \\
CoopGame            & S/N & VideoGame $\sqcap$ $\exists$hasGameMode.CoopMode           & intersectionOf, someValuesFrom \\
% Platform-derived
PCGame              & S/N & VideoGame $\sqcap$ $\exists$availableOn.PCPlatform         & intersectionOf, someValuesFrom \\
ConsoleGame         & S/N & VideoGame $\sqcap$ $\exists$availableOn.ConsolePlatform    & intersectionOf, someValuesFrom \\
MobileGame          & S/N & VideoGame $\sqcap$ $\exists$availableOn.MobilePlatform     & intersectionOf, someValuesFrom \\
NintendoGame        & S/N & VideoGame $\sqcap$ $\exists$availableOn.NintendoPlatform   & intersectionOf, someValuesFrom \\
PlayStationGame     & S/N & VideoGame $\sqcap$ $\exists$availableOn.PlayStationPlatform& intersectionOf, someValuesFrom \\
XboxGame            & S/N & VideoGame $\sqcap$ $\exists$availableOn.XboxPlatform       & intersectionOf, someValuesFrom \\
% Sequel-derived
SequelGame          & S/N & VideoGame $\sqcap$ $\exists$sequelOf.VideoGame             & intersectionOf, someValuesFrom \\
StandaloneGame      & S/N & VideoGame $\sqcap$ $\neg\exists$sequelOf.VideoGame          & intersectionOf, complementOf \\
% Quality-derived
HighlyRatedGame     & S/N & VideoGame $\sqcap$ $\exists$metacriticScore.$\geq_{85}$    & intersectionOf, someValuesFrom, datatype restriction \\
RecentGame          & S/N & VideoGame $\sqcap$ $\exists$releaseDate.$\geq_{2020}$      & intersectionOf, someValuesFrom, datatype restriction \\
% Preesistenti
AwardWinningGame    & S/N & VideoGame $\sqcap$ $\exists$wonAward.Award                 & intersectionOf, someValuesFrom \\
FranchiseGame       & S/N & VideoGame $\sqcap$ $\exists$belongsTo.Franchise            & intersectionOf, someValuesFrom \\
% Cardinality (solo necessarie)
GameWithSingleDeveloper & N & VideoGame $\sqcap$ $\leq1$developedBy                   & subClassOf, maxCardinality \\
ExclusiveRelease    & N & VideoGame $\sqcap$ $\leq1$availableOn                      & subClassOf, maxCardinality \\
\bottomrule
\end{longtable}

\textit{Legenda}: S/N = condizioni sufficienti e necessarie (equivalentClass);
N = solo necessarie (subClassOf).

\subsection{Analisi per Categoria}

\subsubsection{Mode-derived}

Le classi \texttt{MultiplayerGame}, \texttt{SinglePlayerGame} e \texttt{CoopGame}
classificano i giochi in base alla modalit\`a di gioco. Un gioco viene inferito
come \texttt{MultiplayerGame} se esiste almeno un'asserzione \texttt{hasGameMode}
verso un'entit\`a di tipo \texttt{MultiplayerMode}. Si noti che \texttt{SinglePlayerGame}
e \texttt{MultiplayerGame} \textbf{non} sono dichiarate disgiunte: un gioco pu\`o
supportare entrambe le modalit\`a (es. Call of Duty ha sia campagna single-player
che multiplayer).

\textbf{Esempio}: Elden Ring ha \texttt{hasGameMode SinglePlayerMode} e
\texttt{hasGameMode MultiplayerMode} $\rightarrow$ inferito sia come
\texttt{SinglePlayerGame} che come \texttt{MultiplayerGame}.

\subsubsection{Platform-derived}

Sei classi classificano i giochi per piattaforma a tre livelli di granularit\`a:
macrocategoria (\texttt{PCGame}, \texttt{ConsoleGame}, \texttt{MobileGame}) e
famiglia specifica (\texttt{NintendoGame}, \texttt{PlayStationGame}, \texttt{XboxGame}).
L'inferenza sfrutta la gerarchia \texttt{rdfs:subClassOf} delle piattaforme: se
una piattaforma \`e di tipo \texttt{PlayStationPlatform}, il reasoner inferisce
anche \texttt{ConsolePlatform} e \texttt{Platform}.

\textbf{Esempio}: God of War Ragnar\"ok ha \texttt{availableOn} PS5 (tipo
\texttt{PlayStationPlatform}) $\rightarrow$ inferito \texttt{PlayStationGame} e \texttt{ConsoleGame}.

\subsubsection{Sequel-derived}

\texttt{SequelGame} e \texttt{StandaloneGame} formano una partizione completa
dei giochi basata sulla propriet\`a transitiva \texttt{sequelOf}:
\begin{itemize}[nosep]
  \item \texttt{SequelGame} --- $\exists$\texttt{sequelOf.VideoGame}: il gioco ha
    almeno un predecessore diretto;
  \item \texttt{StandaloneGame} --- $\neg\exists$\texttt{sequelOf.VideoGame}: il
    gioco non ha predecessori. Utilizza \texttt{owl:complementOf}, costrutto
    non supportato da OWL 2 RL, e richiede pertanto il reasoner CONSTRUCT
    (cfr.~Sezione~\ref{sec:reasoning}).
\end{itemize}

\textbf{Esempio}: Elden Ring non ha \texttt{sequelOf} $\rightarrow$ \texttt{StandaloneGame};
Dark Souls III ha \texttt{sequelOf} Dark Souls II $\rightarrow$ \texttt{SequelGame}.

\subsubsection{Quality-derived}

\texttt{HighlyRatedGame} e \texttt{RecentGame} utilizzano \textbf{datatype restriction}
con \texttt{xsd:minInclusive}, un costrutto OWL 2 DL non coperto da OWL 2 RL
(cfr.~Lez.~17--18):

\begin{itemize}[nosep]
  \item \texttt{HighlyRatedGame} $\equiv$ VideoGame $\sqcap$ $\exists$\texttt{metacriticScore.}$\geq_{85}$
    --- giochi con Metacritic $\geq$ 85;
  \item \texttt{RecentGame} $\equiv$ VideoGame $\sqcap$ $\exists$\texttt{releaseDate.}$\geq_{2020-01-01}$
    --- giochi rilasciati dal 2020.
\end{itemize}

\textbf{Esempio}: Elden Ring (metacritic 96) $\rightarrow$ \texttt{HighlyRatedGame}.

\subsubsection{Award e Franchise (preesistenti)}

\texttt{AwardWinningGame} e \texttt{FranchiseGame} sono le due classi definite
originali (V1), preservate nell'estensione V2:
\begin{itemize}[nosep]
  \item \texttt{AwardWinningGame}: giochi con almeno un premio (\texttt{wonAward});
  \item \texttt{FranchiseGame}: giochi appartenenti a una serie (\texttt{belongsTo}).
\end{itemize}

\subsubsection{Cardinality (solo necessarie)}

\texttt{GameWithSingleDeveloper} e \texttt{ExclusiveRelease} sono le uniche due
classi definite con \textbf{sole condizioni necessarie} (\texttt{rdfs:subClassOf}
anzich\'e \texttt{owl:equivalentClass}). Questa scelta \`e deliberata:
\begin{itemize}[nosep]
  \item \texttt{GameWithSingleDeveloper}: $\leq 1$ \texttt{developedBy}.
    Condizione sufficiente sarebbe troppo forte (un gioco con esattamente 1
    developer dovrebbe essere \textit{solo} quello);
  \item \texttt{ExclusiveRelease}: $\leq 1$ \texttt{availableOn}.
    Analogo ragionamento: un gioco esclusiva \textit{ha al massimo} 1 piattaforma,
    ma non tutti i giochi con $\leq$ 1 piattaforma sono necessariamente esclusive.
\end{itemize}

Queste classi forniscono vincoli di integrit\`a piuttosto che meccanismi di
classificazione automatica.

\subsection{Costrutti OWL 2 Utilizzati}

La Tabella~\ref{tab:costrutti} riassume i costrutti OWL 2 impiegati nelle classi
definite, con riferimento alle lezioni del corso.

\begin{table}[h]
\centering
\caption{Costrutti OWL 2 nelle classi definite}
\label{tab:costrutti}
\begin{tabular}{lll}
\toprule
\textbf{Costrutto}          & \textbf{Classi}                                   & \textbf{Rif. Lezione} \\
\midrule
\texttt{owl:someValuesFrom}& MultiplayerGame, PCGame, SequelGame, \dots{} (12) & Lez.~15 (esistenziale) \\
\texttt{owl:intersectionOf}& Tutte le 17 definite                             & Lez.~16 (intersezione) \\
\texttt{owl:complementOf}  & StandaloneGame                                   & Lez.~16 (complemento) \\
\texttt{owl:maxCardinality}& GameWithSingleDeveloper, ExclusiveRelease        & Lez.~17--18 \\
\texttt{owl:withRestrictions} & HighlyRatedGame, RecentGame                   & Lez.~17--18 (datatype) \\
\texttt{xsd:minInclusive}  & HighlyRatedGame, RecentGame                      & Lez.~17--18 \\
\bottomrule
\end{tabular}
\end{table}
